\documentclass[12pt]{article}

\usepackage{amsmath,amssymb,booktabs,array,hyperref}
\topmargin -.5cm
\textheight 21cm
\oddsidemargin -.125cm
\textwidth 17cm

\newcommand{\D}{\Delta}
\newcommand{\Tr}{\mathrm{Tr}}
\newcommand{\cO}{\mathcal{O}}
\newcommand{\la}{\langle}
\newcommand{\ra}{\rangle}

\begin{document}
\sloppy

\begin{center}
{\Large\bf Where the Formula for $\D$ Comes From}
\end{center}

\section{Purpose}

This note summarizes the literature origin of two formulas that are easy to conflate but
come from different physics:
\begin{itemize}
\item the exact $\D=4$ of the massless RR axion in $AdS_5\times S^5$;
\item the strong-coupling expansion
\begin{equation}
\D = 2\lambda^{1/4} + (\D_0-4) + \frac{b_1}{\lambda^{1/4}} + \cdots
\end{equation}
for first-excited-string states identified with the Konishi multiplet.
\end{itemize}

Throughout I use
\begin{equation}
R^2/\alpha'=\sqrt{\lambda},\qquad \lambda^{1/4}=R
\end{equation}
in the same conventions as the current SFT notes.

\section{Massless RR Axion}

The RR axion is qualitatively different from a first-massive string mode.  It is already present in
ten-dimensional type IIB supergravity as the massless scalar $C_0$.  Upon reduction on
$AdS_5\times S^5$, the $S^5$ zero harmonic of $C_0$ is therefore a five-dimensional massless scalar.
This is part of the standard Kaluza--Klein spectrum of type IIB supergravity on $S^5$
\cite{KimRomansVN,GunaydinMarcus} and is reviewed from the AdS/CFT side in \cite{AharonyReview}.

For a scalar in $AdS_5$, the standard mass-dimension dictionary is
\begin{equation}
m_5^2 R^2 = \D(\D-4),
\end{equation}
as explained in the original AdS/CFT papers \cite{GKP,Witten}.  For the axion zero mode,
\begin{equation}
m_5^2=0,
\end{equation}
so one gets
\begin{equation}
\D(\D-4)=0.
\end{equation}
The normalizable state in standard quantization is the $\D=4$ branch.  The $\D=0$ branch is the
non-normalizable mode, i.e. changing the boundary value of the source rather than creating a normalizable
bulk state.

So the axion result is
\begin{equation}
\D_{\rm axion}=4.
\end{equation}
There is no $2R$ term because the axion is not a massive string excitation: it is a supergravity zero
mode.  On the CFT side this matches the fact that the axion-dilaton couples to the exactly marginal
complex gauge coupling $\tau$, so the corresponding dimension-four operators are protected
\cite{AharonyReview}.

\subsection*{Direct derivation in the projective $X^a$ coordinates}

The same result can be derived directly in the projective coordinates used in
\texttt{2507.12921}.  For a mode that is trivial on $S^5$, the scalar Laplacian is
\cite{SFT2507}
\begin{equation}
\Box_{\rm AdS}
=
\partial_a\partial^a
+\frac{1}{R^2}\Big(X^aX^b\partial_a\partial_b+5X^a\partial_a\Big)
+O(R^{-4}),
\end{equation}
where $a=0,\ldots,4$.  In the same coordinate frame, the global-time generator is
\begin{equation}
\partial_\tau=X^{-1}\partial_{X^0}-X^0\partial_{X^{-1}},
\qquad
X^{-1}=\sqrt{R^2+X_aX^a},
\end{equation}
and when restricted to functions of the projective coordinates $X^a$ this becomes
\begin{equation}
\partial_\tau
=
\sqrt{R^2+X_aX^a}\,\partial_{X^0}.
\end{equation}
Its large-$R$ expansion is
\begin{equation}
\partial_\tau
=
\left(R+\frac{X_aX^a}{2R}+O(R^{-3})\right)\partial_{X^0}
=
R\partial_{X^0}+\frac{X_aX^a}{2R}\partial_{X^0}+O(R^{-3}),
\end{equation}
as written in \texttt{draftKonishi.tex}.

So a mode of definite global frequency $\omega$ can be written as
\begin{equation}
c_\omega(X)=\Big(\sqrt{R^2+X_aX^a}-iX^0\Big)^{-\omega},
\end{equation}
because if we define
\[
Z(X)\equiv \sqrt{R^2+X_aX^a}-iX^0=X^{-1}-iX^0,
\]
then in global coordinates
\[
X^{-1}=R\sin\tau\cosh\rho,\qquad X^0=R\cos\tau\cosh\rho,
\]
so
\[
Z=R\cosh\rho(\sin\tau-i\cos\tau)=-iR\cosh\rho\,e^{i\tau}.
\]
Therefore $Z$ carries one unit of positive global-time frequency:
\[
\partial_\tau Z=+iZ.
\]
Hence
\begin{equation}
\partial_\tau c_\omega=-i\omega\,c_\omega.
\end{equation}
Indeed,
\[
\partial_\tau c_\omega
=
\partial_\tau(Z^{-\omega})
=
(-\omega)Z^{-\omega-1}(+iZ)
=
-i\omega\,c_\omega.
\]

Now write
\begin{equation}
X_aX^a=-(X^0)^2+r^2,
\qquad
r^2=\sum_{m=1}^4 (X^m)^2.
\end{equation}
Then
\begin{equation}
\sqrt{R^2+X_aX^a}
=
R+\frac{r^2-(X^0)^2}{2R}+O(R^{-3}),
\end{equation}
and hence
\begin{equation}
c_\omega(X)=
R^{-\omega}e^{+i\omega X^0/R}
\left[1-\frac{\omega r^2}{2R^2}+O(R^{-3})\right].
\end{equation}
This expansion is obtained as follows:
\[
Z
=
R\left[
1-\frac{iX^0}{R}+\frac{r^2-(X^0)^2}{2R^2}+O(R^{-3})
\right].
\]
So
\[
Z^{-\omega}
=
R^{-\omega}(1+u)^{-\omega},
\qquad
u=-\frac{iX^0}{R}+\frac{r^2-(X^0)^2}{2R^2}+O(R^{-3}).
\]
Using
\[
(1+u)^{-\omega}=1-\omega u+\frac{\omega(\omega+1)}{2}u^2+O(u^3),
\qquad
u^2=-\frac{(X^0)^2}{R^2}+O(R^{-3}),
\]
one gets
\[
Z^{-\omega}
=
R^{-\omega}\left[
1+\frac{i\omega X^0}{R}
-\frac{\omega r^2}{2R^2}
-\frac{\omega^2(X^0)^2}{2R^2}
+O(R^{-3})
\right].
\]
The $X^0$-dependent terms resum into
\[
e^{+i\omega X^0/R}
=
1+\frac{i\omega X^0}{R}-\frac{\omega^2(X^0)^2}{2R^2}+O(R^{-3}),
\]
which gives the displayed form.

Acting with the Laplacian gives
\begin{equation}
\Box_{\rm AdS}c_\omega
=
\frac{\omega(\omega-4)}{R^2}c_\omega+O(R^{-3}),
\end{equation}
and the derivation is straightforward.  Write
\[
c_\omega(X)=R^{-\omega}e^{+i\omega X^0/R}f(r),
\qquad
f(r)=1-\frac{\omega r^2}{2R^2}+O(R^{-3}).
\]
To this order, the explicit curvature correction in $\Box_{\rm AdS}$ does not contribute:
\[
\frac{1}{R^2}\Big(X^aX^b\partial_a\partial_b+5X^a\partial_a\Big)c_\omega=O(R^{-3}),
\]
because every derivative acting on the leading mode brings at least one extra factor of $1/R$.
So to order $R^{-2}$,
\[
\Box_{\rm AdS}c_\omega=\partial_a\partial^a c_\omega+O(R^{-3})
=
\left(-\partial_0^2+\sum_{m=1}^4\partial_m^2\right)c_\omega+O(R^{-3}).
\]
For the time derivative,
\[
-\partial_0^2 c_\omega
=
\frac{\omega^2}{R^2}c_\omega+O(R^{-3}),
\]
since only the exponential depends on $X^0$ at this order.  For the spatial derivatives,
\[
\sum_{m=1}^4\partial_m^2 c_\omega
=
R^{-\omega}e^{+i\omega X^0/R}\sum_{m=1}^4\partial_m^2
\left(1-\frac{\omega r^2}{2R^2}\right)+O(R^{-3}).
\]
Using
\[
\sum_{m=1}^4\partial_m^2 r^2 = 2\times 4 = 8,
\]
this becomes
\[
\sum_{m=1}^4\partial_m^2 c_\omega
=
-\frac{4\omega}{R^2}c_\omega+O(R^{-3}).
\]
Adding time and space pieces,
\[
\Box_{\rm AdS}c_\omega
=
\left(\frac{\omega^2}{R^2}-\frac{4\omega}{R^2}\right)c_\omega+O(R^{-3})
=
\frac{\omega(\omega-4)}{R^2}c_\omega+O(R^{-3}),
\]
which is exactly the displayed equation.

so the wave equation $\Box_{\rm AdS}c_\omega=0$ implies
\begin{equation}
\omega(\omega-4)=0.
\end{equation}
This is just the condition that the coefficient of the first nontrivial term in the large-$R$ expansion
vanish.  So there are two branches,
\[
\omega=0,\qquad \omega=4.
\]
The $\omega=0$ solution is the non-normalizable/source branch, while the normalizable massless scalar is
the $\omega=4$ branch.

The normalizable branch is therefore
\begin{equation}
\omega=4.
\end{equation}
Thus the projective-coordinate derivation reproduces the standard massless-scalar result
\begin{equation}
\D=4.
\end{equation}

\section{The Same Calculation for a Massive Scalar}

It is useful to repeat exactly the same projective-coordinate calculation for a single massive scalar with
trivial $S^5$ dependence.  The equation of motion is now
\begin{equation}
(\Box_{\rm AdS}-m^2)\phi=0.
\end{equation}

Take the same definite-frequency ansatz
\begin{equation}
\phi_\omega(X)=\Big(\sqrt{R^2+X_aX^a}-iX^0\Big)^{-\omega}.
\end{equation}
Nothing changes about the time dependence: it is still true that
\begin{equation}
\partial_\tau\phi_\omega=-i\omega\,\phi_\omega.
\end{equation}
And the large-$R$ expansion is the same as before,
\begin{equation}
\phi_\omega(X)=
R^{-\omega}e^{+i\omega X^0/R}
\left[1-\frac{\omega r^2}{2R^2}+O(R^{-3})\right].
\end{equation}

Likewise, the Laplacian acts in the same way:
\begin{equation}
\Box_{\rm AdS}\phi_\omega
=
\frac{\omega(\omega-4)}{R^2}\phi_\omega+O(R^{-3}).
\end{equation}
Plugging this into the massive wave equation gives
\begin{equation}
\frac{\omega(\omega-4)}{R^2}\phi_\omega-m^2\phi_\omega+O(R^{-3})=0,
\end{equation}
so at leading nontrivial order one gets
\begin{equation}
\omega(\omega-4)=m^2R^2.
\end{equation}
This is the standard AdS$_5$ massive-scalar relation.  Solving the quadratic equation gives
\begin{equation}
\omega_\pm=2\pm \sqrt{4+m^2R^2}.
\end{equation}
The normalizable branch is therefore
\begin{equation}
\omega=\D=2+\sqrt{4+m^2R^2}.
\end{equation}

For large mass, one may expand
\begin{equation}
\D
=
2+\sqrt{4+m^2R^2}
=
mR+2+\frac{2}{mR}+O((mR)^{-3}).
\end{equation}
Thus a \emph{single} massive scalar field in AdS has a universal constant offset $+2$ at large mass.
It does not produce a state-dependent offset of the form $\D_0-4$.

This is precisely why a naive one-field massive-scalar wave equation is not enough to derive the Konishi
formula.  Such a wave equation knows only about the bulk mass $m$, whereas the Konishi constant term
depends on the position of the state inside the full supermultiplet.

\subsection*{Universal time-dependence factor}

There is nevertheless one useful universal lesson from the scalar calculation.  The truly universal
statement is not that every mode is literally equal to $Z^{-\omega}$ times a polynomial.  Rather, it is
that on the static $AdS_5\times S^5$ background every linearized field may be decomposed into eigenmodes
of the global-time Killing vector $\partial_\tau$.  For any such mode one can therefore factor out
\begin{equation}
e^{-i\omega\tau}.
\end{equation}
This holds for any bosonic or fermionic fluctuation, provided one works in a $\tau$-independent basis for
the tensor or spinor indices.  The reason is simply that the linearized kinetic operators commute with
the isometry generated by $\partial_\tau$, so one may choose simultaneous eigenmodes of the equations of
motion and of $i\partial_\tau$.  Sphere quantum numbers do not change this statement: they live in the
$S^5$ factor and are spectators as far as global AdS time is concerned.

In projective coordinates it is convenient to trade the explicit factor $e^{-i\omega\tau(X)}$ for a power
of the combination
\begin{equation}
Z\equiv X^{-1}-iX^0=\sqrt{R^2+X_aX^a}-iX^0.
\end{equation}
Since
\begin{equation}
\partial_\tau Z=+iZ,
\end{equation}
the factor $Z$ has unit $\tau$-charge.  Moreover,
\begin{equation}
X^{-1}-iX^0=-i\sqrt{R^2+r^2}\,e^{i\tau},
\end{equation}
so
\begin{equation}
e^{-i\omega\tau}
=
(-i)^\omega \big(R^2+r^2\big)^{-\omega/2} Z^{-\omega}.
\end{equation}
Therefore any definite-frequency mode can be rewritten as
\begin{equation}
\Phi_\omega(X,U)=e^{-i\omega\tau(X)}\,\Psi(X^m,U)
=
Z^{-\omega}\,\widetilde\Psi(X^m,U),
\end{equation}
where
\begin{equation}
\widetilde\Psi(X^m,U)=(-i)^\omega \big(R^2+r^2\big)^{\omega/2}\Psi(X^m,U),
\end{equation}
and $U$ denotes the sphere variables.  Since $r^2$ and the sphere coordinates are $\tau$-independent,
this is just a reshuffling of the radial profile.  In this precise sense the $Z^{-\omega}$ factor is
universal in projective coordinates: it is another way of writing the universal $e^{-i\omega\tau}$ time
dependence.

Near the center of AdS this becomes
\begin{equation}
Z^{-\omega}\sim (-i)^{-\omega}R^{-\omega}e^{-i\omega\tau}
\sim (-i)^{-\omega}R^{-\omega}e^{-i\omega X^0/R}.
\end{equation}

What is not universal is the remaining profile $\widetilde\Psi$.  For a generic mode it is not a simple
polynomial; it is determined by the radial equations, tensor constraints, and possible component mixing.
So for a more general state, the natural projective-coordinate ansatz is
\begin{equation}
\Phi_\omega(X,U)=Z^{-\omega}\times \text{(remaining $\tau$-independent structure)}.
\end{equation}
For lowest-weight states that remaining structure often starts with a simple tensor polynomial.  For
example, an AdS highest-weight state carrying spin in a fixed plane is expected to look schematically like
\begin{equation}
\Phi_\omega \sim Z^{-\omega-s}(X^+)^s \times \text{(rest of the state)}.
\end{equation}
Equivalently, with a null polarization vector $n$, one may write the kinematic factor as
\begin{equation}
\Phi_\omega \sim Z^{-\omega-s}(n\cdot X)^s.
\end{equation}

So the clean statement is:
\begin{itemize}
\item $e^{-i\omega\tau}$ is universal for any definite-frequency mode on $AdS_5\times S^5$;
\item $Z^{-\omega}$ is the projective-coordinate form of that universal time factor, after absorbing a
      $\tau$-independent radial factor into the rest of the wavefunction;
\item the actual relation between $\omega$ and the physical state data, as well as the mixing between
      different components, is dynamical and cannot be read off from this factor alone.
\end{itemize}

\subsection*{Starting instead from \texorpdfstring{$e^{-i\omega\tau}f(\rho)$}{exp(-i omega tau) f(rho)}}

If one prefers, one can start from the more primitive global-coordinate ansatz
\begin{equation}
\phi(\tau,\rho)=e^{-i\omega\tau}f(\rho)
\end{equation}
for an $S^3$-invariant scalar mode.  In global $AdS_5$ with metric
\begin{equation}
ds^2=R^2\left(-\cosh^2\rho\,d\tau^2+d\rho^2+\sinh^2\rho\,d\Omega_3^2\right),
\end{equation}
the massive scalar equation $(\Box_{\rm AdS}-m^2)\phi=0$ becomes
\begin{equation}
\frac{1}{\sinh^3\rho\,\cosh\rho}\partial_\rho
\Big(\sinh^3\rho\,\cosh\rho\,\partial_\rho f\Big)
\,+\,\frac{\omega^2}{\cosh^2\rho}f
\,-\,m^2R^2 f
=
0.
\end{equation}
Now try
\begin{equation}
f(\rho)=(\cosh\rho)^{-\omega}.
\end{equation}
Then
\[
\partial_\rho f=-\omega\tanh\rho\,f,
\]
so
\[
\sinh^3\rho\,\cosh\rho\,\partial_\rho f
=
-\omega\sinh^4\rho\,f.
\]
Differentiating once more,
\[
\partial_\rho\Big(\sinh^3\rho\,\cosh\rho\,\partial_\rho f\Big)
=
-\omega\Big(4\sinh^3\rho\cosh\rho\,f+\sinh^4\rho\,\partial_\rho f\Big),
\]
which gives
\[
\frac{1}{\sinh^3\rho\,\cosh\rho}\partial_\rho
\Big(\sinh^3\rho\,\cosh\rho\,\partial_\rho f\Big)
=
-4\omega f+\omega^2\tanh^2\rho\,f.
\]
Adding the time-derivative term,
\[
-4\omega f+\omega^2\tanh^2\rho\,f+\frac{\omega^2}{\cosh^2\rho}f
=
\omega(\omega-4)f,
\]
since $\tanh^2\rho+\cosh^{-2}\rho=1$.  So the equation is solved provided
\begin{equation}
\omega(\omega-4)=m^2R^2.
\end{equation}

Now rewrite the solution in projective coordinates.  From
\begin{equation}
X^{-1}=R\sin\tau\cosh\rho,\qquad X^0=R\cos\tau\cosh\rho,
\end{equation}
one gets
\begin{equation}
X^{-1}-iX^0=-iR\cosh\rho\,e^{i\tau}.
\end{equation}
Therefore
\begin{equation}
e^{-i\omega\tau}(\cosh\rho)^{-\omega}
=
(-iR)^\omega\,(X^{-1}-iX^0)^{-\omega},
\end{equation}
up to an overall constant phase.  So the projective-coordinate factor $Z^{-\omega}$ is not ad hoc:
it is exactly the global-coordinate ansatz $e^{-i\omega\tau}f(\rho)$ with the lowest-weight radial
profile $f(\rho)=(\cosh\rho)^{-\omega}$ rewritten in terms of the projective coordinates.

\subsection*{The same ansatz, but entirely in projective coordinates}

If one wants to avoid $(\tau,\rho)$ as independent coordinates altogether, one can still start from the
same idea directly in the projective variables.  The global time is a function of the projective
coordinates:
\begin{equation}
\tan\tau(X)=\frac{X^{-1}}{X^0}
=
\frac{\sqrt{R^2+X_aX^a}}{X^0},
\end{equation}
or, more usefully,
\begin{equation}
e^{-i\tau(X)}=\frac{X^0-iX^{-1}}{\sqrt{(X^{-1})^2+(X^0)^2}}
=
\frac{X^0-i\sqrt{R^2+X_aX^a}}{\sqrt{R^2+r^2}},
\end{equation}
where
\begin{equation}
r^2=\sum_{m=1}^4 (X^m)^2,
\qquad
(X^{-1})^2+(X^0)^2=R^2+r^2.
\end{equation}

The exact induced AdS$_5$ metric in these coordinates is
\begin{equation}
G_{ab}=\eta_{ab}-\frac{X_aX_b}{R^2+X_cX^c},
\end{equation}
with inverse and determinant
\begin{equation}
G^{ab}=\eta^{ab}+\frac{X^aX^b}{R^2},
\qquad
\sqrt{-G}=\frac{R}{\sqrt{R^2+X_cX^c}}.
\end{equation}
Therefore the exact scalar Laplacian is
\begin{equation}
\Box_{\rm AdS}
=
\frac{1}{\sqrt{-G}}\partial_a\!\left(\sqrt{-G}\,G^{ab}\partial_b\right)
=
\left(\eta^{ab}+\frac{X^aX^b}{R^2}\right)\partial_a\partial_b
+\frac{5X^a}{R^2}\partial_a .
\end{equation}
Expanding this exact expression reproduces the earlier large-$R$ form
\[
\Box_{\rm AdS}
=
\partial_a\partial^a
+\frac{1}{R^2}\Big(X^aX^b\partial_a\partial_b+5X^a\partial_a\Big)
+O(R^{-4}).
\]

Now keep the explicit time dependence and make a genuinely radial ansatz
\begin{equation}
\phi(X)=e^{-i\omega\tau(X)}\,f(s),
\qquad
s\equiv \frac{r^2}{R^2}.
\end{equation}
The needed derivatives in $X$ coordinates are straightforward.  Since
\[
s=\frac{1}{R^2}\sum_{m=1}^4 (X^m)^2,
\]
one has
\begin{equation}
\partial_0 s=0,
\qquad
\partial_m s=\frac{2X_m}{R^2}.
\end{equation}
For the global time, differentiating
\[
\tan\tau(X)=\frac{X^{-1}}{X^0},
\qquad
X^{-1}=\sqrt{R^2+X_aX^a},
\]
gives
\begin{equation}
\partial_0\tau=-\frac{1}{X^{-1}},
\qquad
\partial_m\tau=\frac{X^0X_m}{X^{-1}(R^2+r^2)}.
\end{equation}

Now contract these derivatives with the exact inverse metric.  A direct computation gives
\begin{equation}
G^{ab}\partial_a\tau\,\partial_b\tau
=
-\frac{1}{R^2+r^2}
=
-\frac{1}{R^2(1+s)},
\end{equation}
\begin{equation}
G^{ab}\partial_a\tau\,\partial_b s=0,
\end{equation}
\begin{equation}
G^{ab}\partial_a s\,\partial_b s
=
\frac{4r^2}{R^4}+\frac{4r^4}{R^6}
=
\frac{4s(1+s)}{R^2},
\end{equation}
and
\begin{equation}
\Box_{\rm AdS}\tau=0,
\qquad
\Box_{\rm AdS}s=\frac{4(2+3s)}{R^2}.
\end{equation}

With these ingredients, the chain rule gives
\begin{align}
\Box_{\rm AdS}\phi
&=
e^{-i\omega\tau}\Big[
-\omega^2\,G^{ab}\partial_a\tau\,\partial_b\tau\, f
-2i\omega\,G^{ab}\partial_a\tau\,\partial_b s\, f'
-i\omega\,(\Box_{\rm AdS}\tau)\,f
\notag\\
&\hspace{3.8cm}
+\,G^{ab}\partial_a s\,\partial_b s\, f''
+(\Box_{\rm AdS}s)\,f'
\Big].
\end{align}
Substituting the four scalar identities above, the mixed term drops out and one obtains the exact ODE
\begin{equation}
4s(1+s)f''(s)+4(2+3s)f'(s)+\frac{\omega^2}{1+s}f(s)=0.
\end{equation}

This is the differential equation for the radial profile, derived directly from the $X$-coordinate
metric and Laplacian.  This is the reusable step: for a more complicated state, such as a Konishi
fluctuation with additional tensor structure and component mixing, one would first derive the analogue
of this radial operator and only then study its coupled system.  Only now does one make a
lowest-weight ansatz for the scalar solution.  Trying
\begin{equation}
f(s)=(1+s)^{-\beta},
\end{equation}
one gets
\begin{equation}
\omega^2-8\beta=0,
\qquad
4\beta^2-8\beta=0.
\end{equation}
So
\begin{equation}
\beta=0\ \text{or}\ 2.
\end{equation}
The $\beta=0$ branch is the non-normalizable constant solution, which forces $\omega=0$.  The
normalizable lowest-weight branch is
\begin{equation}
\beta=2,
\qquad
\omega=4.
\end{equation}

Therefore
\begin{equation}
\phi(X)=e^{-i4\tau(X)}\left(1+\frac{r^2}{R^2}\right)^{-2}.
\end{equation}
Using
\begin{equation}
X^{-1}-iX^0=-i\sqrt{R^2+r^2}\,e^{i\tau},
\end{equation}
this may be rewritten as
\begin{equation}
\phi(X)\propto \left(X^{-1}-iX^0\right)^{-4}.
\end{equation}
So the compact projective factor is not a guess: it is the solution of the exact $X$-coordinate radial
equation for the lowest normalizable mode.

\subsection*{Large-$R$ series solution of the radial equation}

One can also solve the same radial equation locally, as a large-$R$ expansion, without guessing the
closed-form answer.  Write
\begin{equation}
f(s)=\sum_{n=0}^\infty a_n s^n,
\qquad
s=\frac{r^2}{R^2}.
\end{equation}
Then
\begin{equation}
f'(s)=\sum_{n=1}^\infty n a_n s^{n-1},
\qquad
f''(s)=\sum_{n=2}^\infty n(n-1)a_n s^{n-2},
\end{equation}
and the exact radial equation
\begin{equation}
4s(1+s)f''(s)+4(2+3s)f'(s)+\frac{\omega^2}{1+s}f(s)=0
\end{equation}
may be expanded around $s=0$ using
\begin{equation}
\frac{1}{1+s}=1-s+s^2-\cdots.
\end{equation}
Keeping the first few powers of $s$ gives
\begin{align}
0
&=
8a_1+\omega^2 a_0 \notag\\
&\quad
+\Big(24a_2+12a_1+\omega^2(a_1-a_0)\Big)s \notag\\
&\quad
+\Big(48a_3+32a_2+\omega^2(a_2-a_1+a_0)\Big)s^2
+O(s^3).
\end{align}
So the first coefficients are fixed recursively in terms of $a_0$ and $\omega$:
\begin{equation}
a_1=-\frac{\omega^2}{8}a_0,
\end{equation}
\begin{equation}
a_2=\frac{\omega^2(\omega^2+20)}{192}a_0,
\end{equation}
\begin{equation}
a_3=-\frac{\omega^2(\omega^4+44\omega^2+384)}{9216}a_0.
\end{equation}
Therefore the local large-$R$ solution is
\begin{equation}
f\!\left(\frac{r^2}{R^2}\right)
=
a_0\left[
1
-\frac{\omega^2}{8}\frac{r^2}{R^2}
+\frac{\omega^2(\omega^2+20)}{192}\frac{r^4}{R^4}
+O(R^{-6})
\right].
\end{equation}
Equivalently,
\begin{equation}
\phi(X)
=
e^{-i\omega\tau(X)}\,a_0\left[
1
-\frac{\omega^2}{8}\frac{r^2}{R^2}
+\frac{\omega^2(\omega^2+20)}{192}\frac{r^4}{R^4}
+O(R^{-6})
\right].
\end{equation}

This expansion does \emph{not} quantize $\omega$.  For any chosen $\omega$, there is a regular local
solution near the origin, with coefficients determined recursively by the equation.  The value
$\omega=4$ is fixed only after imposing the global AdS boundary condition: one must select the
normalizable lowest-weight mode on the full spacetime.  For that branch, the exact solution found above
is
\begin{equation}
f(s)=(1+s)^{-2}=1-2s+3s^2-4s^3+\cdots,
\end{equation}
which is recovered from the recursion precisely when $\omega=4$.

\paragraph{Comment on the general mechanism.}
This scalar example makes clear what the large-radius expansion can and cannot do.  The inner expansion
at fixed projective coordinates $X^a$ probes only the region $s=r^2/R^2\ll 1$, so it determines a local
regular solution but does not quantize $\omega$.  The quantization of $\omega$ is a genuine boundary-value
effect: one must continue the solution to the AdS boundary, $s\to\infty$, and require the absence of the
non-normalizable branch.  For a scalar of mass $m$ the large-$s$ equation gives
\begin{equation}
f(s)\sim A\,s^{-\D/2}+B\,s^{-(4-\D)/2},
\qquad
\D=2+\sqrt{4+m^2R^2},
\end{equation}
and the AdS boundary condition is the vanishing of the coefficient of the bad branch.  For generic
Konishi fluctuations the same structure should hold at the level of reduced radial amplitudes in a fixed
harmonic basis: after extracting the universal $e^{-i\omega\tau}$ factor, one obtains a coupled radial
system, and the allowed values of $\omega$ are those for which the source/non-normalizable asymptotics are
absent.  What the local large-$R$ expansion provides is the inner data for such a matching problem, not by
itself the full frequency quantization.

\subsection*{Generic boundary condition for a primary mode}

For a generic mode on $AdS_5\times S^5$, including spinning modes, the right boundary condition is not
``primary'' by itself, but rather:
\begin{equation}
\text{regular at the center}
\quad + \quad
\text{source-free at the AdS boundary}.
\end{equation}
The distinction between primary and descendant is then made \emph{within} the space of normalizable
solutions: the primary is the lowest-frequency state in the corresponding $SO(2,4)$ module.

More concretely, one should first decompose the fluctuation into a fixed tangent-frame
$SO(4)\times SO(5)$ harmonic basis and extract the universal factor $e^{-i\omega\tau}$.  The remaining
radial amplitudes satisfy a coupled ordinary differential system in the AdS radial variable.  Near the
AdS boundary there are always two types of asymptotic behavior for these reduced amplitudes:
\begin{equation}
\psi_{\rm red}\sim \psi_{\rm source}+\psi_{\rm vev},
\end{equation}
namely a non-normalizable/source branch and a normalizable/state branch.  For a scalar these reduce to
\begin{equation}
\psi_{\rm red}\sim A\,z^{4-\D}+B\,z^\D,
\qquad z\to 0,
\end{equation}
or equivalently
\begin{equation}
\psi_{\rm red}\sim A\,e^{-(4-\D)\rho}+B\,e^{-\D\rho},
\qquad \rho\to\infty.
\end{equation}
For higher spin the raw coordinate components need not follow this literal power law because the vielbein,
constraints, and gauge conditions contribute extra powers.  The universal statement is only for the
reduced tangent-frame amplitudes after the kinematic polarization factors are stripped off.

Thus the generic AdS boundary condition for the state is
\begin{equation}
\psi_{\rm source}=0.
\end{equation}
This condition quantizes $\omega$.  Once one has the discrete normalizable spectrum in a given
representation, the primary is the solution with the smallest allowed global energy, while the
descendants are obtained by acting with raising operators and have larger $\omega$.

\section{First Excited String Level}

The Konishi formula starts from a completely different input.  One is no longer talking about a
massless supergravity fluctuation, but about a short string on the first excited closed-string level.
So there is no derivation parallel to the axion case where one solves a local AdS wave equation and reads
off $\D$ from near-boundary falloff.  For the massive string state, the literature derivation is split:
\begin{itemize}
\item the universal leading $2\lambda^{1/4}$ is derived from the first-massive string mass;
\item the constant term is fixed by the Konishi supermultiplet assignment and is tied to zero modes.
\end{itemize}

\subsection*{Leading term}

The closed-string flat-space mass-shell condition is
\begin{equation}
\alpha' m^2 = 4(N-1),
\end{equation}
so for the first excited level, $N=2$ and thus
\begin{equation}
m^2_{\rm 1st\ excited}=\frac{4}{\alpha'}
\end{equation}
for the first massive closed-string level.  For a heavy state in AdS one then expects
\begin{equation}
\D \sim E \sim mR = \frac{2R}{\sqrt{\alpha'}} = 2\lambda^{1/4}.
\end{equation}
This is the same qualitative strong-coupling scaling emphasized already in the original AdS/CFT
literature for massive string states \cite{GKP,AharonyReview}.

\subsection*{Constant term}

The more precise semiclassical short-string analysis of Roiban and Tseytlin writes the first-excited
level energy as \cite{RT2011}
\begin{equation}
E = 2\lambda^{1/4}+b_0+\frac{b_1}{\lambda^{1/4}}+\frac{b_2}{\lambda^{1/2}}+\cdots .
\end{equation}
They also state that for the Konishi multiplet
\begin{equation}
b_0+4=\D_0,
\end{equation}
so
\begin{equation}
b_0=\D_0-4.
\end{equation}

The leading coefficient $2$ and the constant term $b_0$ have different origins:
\begin{itemize}
\item $2\lambda^{1/4}$ is the flat-space first-massive string energy written in AdS units;
\item $b_0=\D_0-4$ is the multiplet-position or zero-mode offset that distinguishes different states
inside the same long Konishi multiplet.
\end{itemize}

Roiban and Tseytlin are explicit that the constant term is not derived in the same direct way as the
leading classical energy.  They say that a systematic derivation of this contribution requires
zero-mode superparticle quantization \cite{RT2011}.  In other words, the literature does \emph{not}
say that $(\D_0-4)$ is another flat-space oscillator mass.  It is an $O(1)$ term associated with the
AdS supermultiplet structure of the state.

So the precise status of the formula is:
\begin{equation}
\D = 2\lambda^{1/4}+O(1)
\end{equation}
is directly derived from the first-massive string mass, while
\begin{equation}
\D = 2\lambda^{1/4}+(\D_0-4)+\frac{b_1}{\lambda^{1/4}}+\cdots
\end{equation}
is the Konishi multiplet formula used in the short-string literature, with the constant term coming from
the supermultiplet and zero-mode sector rather than from a solved local field equation.

For example, on the first excited level:
\begin{equation}
\D_0=2 \Rightarrow \D = 2R-2+\cdots,\qquad
\D_0=4 \Rightarrow \D = 2R+\cdots,\qquad
\D_0=6 \Rightarrow \D = 2R+2+\cdots .
\end{equation}

The one-loop and pure-spinor literature then checks parts of this strong-coupling picture from
independent angles.  In particular, Roiban and Tseytlin compute the semiclassical correction to $b_1$
\cite{RT2009,RT2011}, while Vallilo and Mazzucato give an independent pure-spinor analysis of the
Konishi multiplet at strong coupling \cite{ValliloMazzucato}.

\section{Comparison Table}

\begin{center}
\small
\renewcommand{\arraystretch}{1.25}
\begin{tabular}{>{\raggedright\arraybackslash}p{2.6cm} >{\raggedright\arraybackslash}p{3.7cm} >{\raggedright\arraybackslash}p{3.4cm} >{\raggedright\arraybackslash}p{4.2cm}}
\toprule
state & bulk origin & formula for $\D$ & literature origin \\
\midrule
massless RR axion &
supergravity zero mode of $C_0$ on $S^5$ &
$m_5^2=0 \Rightarrow \D(\D-4)=0 \Rightarrow \D=4$ for the normalizable branch &
IIB KK spectrum on $S^5$ plus the AdS scalar dictionary \cite{KimRomansVN,GunaydinMarcus,GKP,Witten,AharonyReview} \\
\midrule
first excited string state in Konishi &
massive closed-string mode at the first excited level &
$\D=2\lambda^{1/4}+(\D_0-4)+b_1/\lambda^{1/4}+\cdots$ &
short-string semiclassics and Konishi multiplet assignment \cite{RT2009,RT2011}; pure-spinor cross-check \cite{ValliloMazzucato} \\
\bottomrule
\end{tabular}
\end{center}

\section{What This Means for the Current Puzzle}

The literature distinction is sharp.

For the axion, the statement $\D=4$ is a direct consequence of being a massless supergravity scalar in
$AdS_5$, together with the standard choice of normalizable branch.  This is why the axion is a clean
normalization condition for the Hamiltonian: there is no competing $2R$ piece to begin with.

For first-massive Konishi states, the story is different.  The leading $2R$ comes from the universal
flat-space string mass, while the constant piece $\D_0-4$ is a separate $O(1)$ effect tied to the
AdS supermultiplet structure.  The literature supports the existence of this term, but also says that a
systematic worldsheet derivation of it requires the zero modes.  That is exactly the point where the
current SFT contradiction is sharpest: the representation-theory scan says some $\D_0=6$ states have no
obvious first-order mixing partner, while the Konishi formula still demands $b_0=2$.

\appendix

\section{What Is Actually Derived}

The current note uses the following hierarchy of claims.

\paragraph{A. Axion.}
The precise chain is:
\begin{equation}
C_0\ \text{is a massless IIB field}
\ \Longrightarrow \
\text{$S^5$ zero harmonic gives }m_5^2=0
\ \Longrightarrow \
\D(\D-4)=0
\ \Longrightarrow \
\D=4
\end{equation}
for the normalizable state.  This is a direct supergravity plus AdS/CFT derivation.

\paragraph{B. First excited level: leading term.}
The precise chain is:
\begin{equation}
m^2_{\rm 1st\ excited}=\frac{4}{\alpha'}
\ \Longrightarrow \
\D\sim mR=2\lambda^{1/4}.
\end{equation}
This is the universal large-$\lambda$ piece for first-massive string states.

\paragraph{C. First excited level: constant term.}
The literature formula
\begin{equation}
b_0=\D_0-4
\end{equation}
is part of the Konishi multiplet assignment in the short-string analysis, but Roiban and Tseytlin
explicitly note that a systematic derivation of this contribution requires zero-mode superparticle
quantization.  So the existence of the term is part of the accepted strong-coupling picture, but its
worldsheet derivation is subtler than the derivation of the leading $2\lambda^{1/4}$.

\paragraph{D. Why the axion does not help derive $b_0$.}
The axion and the first excited level probe different sectors:
\begin{itemize}
\item the axion is a protected supergravity mode with $\D=O(1)$ and a standard AdS wave equation;
\item the Konishi state is a non-BPS short string with $\D\sim 2R$ and no comparable one-field wave
equation derivation.
\end{itemize}
So the axion fixes the normalization of the Hamiltonian on the massless sector, but by itself it does
not explain where the Konishi constant term $\D_0-4$ comes from.

\begin{thebibliography}{99}

\bibitem{GKP}
S.~S.~Gubser, I.~R.~Klebanov and A.~M.~Polyakov,
``Gauge theory correlators from non-critical string theory,''
Phys.\ Lett.\ B {\bf 428} (1998) 105,
\href{https://arxiv.org/abs/hep-th/9802109}{hep-th/9802109}.

\bibitem{Witten}
E.~Witten,
``Anti-de Sitter space and holography,''
Adv.\ Theor.\ Math.\ Phys.\ {\bf 2} (1998) 253,
\href{https://arxiv.org/abs/hep-th/9802150}{hep-th/9802150}.

\bibitem{KimRomansVN}
H.~J.~Kim, L.~J.~Romans and P.~van Nieuwenhuizen,
``The mass spectrum of chiral $N=2$, $D=10$ supergravity on $S^5$,''
Phys.\ Rev.\ D {\bf 32} (1985) 389.

\bibitem{GunaydinMarcus}
M.~Gunaydin and N.~Marcus,
``The spectrum of the $S^5$ compactification of the chiral $N=2$, $D=10$ supergravity and the
unitary supermultiplets of $U(2,2/4)$,''
Class.\ Quant.\ Grav.\ {\bf 2} (1985) L11.

\bibitem{AharonyReview}
O.~Aharony, S.~S.~Gubser, J.~M.~Maldacena, H.~Ooguri and Y.~Oz,
``Large $N$ field theories, string theory and gravity,''
Phys.\ Rept.\ {\bf 323} (2000) 183,
\href{https://arxiv.org/abs/hep-th/9905111}{hep-th/9905111}.

\bibitem{RT2009}
R.~Roiban and A.~A.~Tseytlin,
``Quantum strings in $AdS_5\times S^5$: strong-coupling corrections to dimension of Konishi operator,''
\href{https://arxiv.org/abs/0906.4294}{arXiv:0906.4294}.

\bibitem{RT2011}
R.~Roiban and A.~A.~Tseytlin,
``Semiclassical string computation of strong-coupling corrections to dimensions of operators in Konishi multiplet,''
Nucl.\ Phys.\ B {\bf 848} (2011) 251,
\href{https://arxiv.org/abs/1102.1209}{arXiv:1102.1209}.

\bibitem{ValliloMazzucato}
B.~C.~Vallilo and L.~Mazzucato,
``The Konishi multiplet at strong coupling,''
JHEP {\bf 12} (2011) 029,
\href{https://arxiv.org/abs/1102.1219}{arXiv:1102.1219}.

\bibitem{SFT2507}
R.~Pius, B.~T.~Stefa\'nski and A.~A.~Tseytlin,
``Type IIB string field theory on $AdS_5\times S^5$ at large radius,''
\href{https://arxiv.org/abs/2507.12921}{arXiv:2507.12921}.

\end{thebibliography}

\end{document}
